The reformulation of Section~\ref{sec:method} is stated in terms of algebra and
memory traffic, so it carries over to any accelerator whose cost is dominated by
traffic to a large, slow memory. We re-derived the same two kernels twice more
as a test of that claim. Table~\ref{tab:backends} collects the results.

% Generated by scripts/seed_from_report.py from the authors' original
% measurement campaign. Re-running scripts/bench.py followed by
% scripts/make_tables.py regenerates this file from fresh measurements.
\begin{table}[t]
\centering
\caption{Speedup over the reference implementation for all three realisations of the same reformulation, homogeneous protocol ($k{=}3$). The Metal figures are measured against PyTorch's MPS backend, which is markedly less mature than cuDNN; they demonstrate portability of the formulation rather than a like-for-like vendor comparison.}
\label{tab:backends}
\begin{tabular}{lllccccc}
\toprule
\multirow{2}{*}{Backend} & \multirow{2}{*}{Hardware} & \multirow{2}{*}{Precision} & \multicolumn{5}{c}{Decomposition levels $L$} \\
\cmidrule(lr){4-8}
 & & & 1 & 2 & 3 & 4 & 5 \\
\midrule
CUDA & RTX A6000 & \texttt{fp32} & 2.24$\times$ & 2.72$\times$ & 2.82$\times$ & 2.85$\times$ & 2.87$\times$ \\
 &  & \texttt{fp16} & 2.99$\times$ & 3.65$\times$ & 3.79$\times$ & 3.83$\times$ & 3.84$\times$ \\
\addlinespace[2pt]
Triton & RTX A6000 & \texttt{fp32} & 2.71$\times$ & 3.09$\times$ & 3.03$\times$ & 2.90$\times$ & 2.81$\times$ \\
 &  & \texttt{fp16} & 2.94$\times$ & 3.16$\times$ & 3.00$\times$ & 2.81$\times$ & 2.69$\times$ \\
\addlinespace[2pt]
Metal & M3 Pro & \texttt{fp32} & 2.05$\times$ & 2.34$\times$ & 2.40$\times$ & 2.37$\times$ & 2.31$\times$ \\
 &  & \texttt{fp16} & 2.05$\times$ & 2.13$\times$ & 2.20$\times$ & 2.20$\times$ & 2.15$\times$ \\
\addlinespace[2pt]
\bottomrule
\end{tabular}
\end{table}


\subsection{Triton}

The Triton realisation \citep{tillet2019triton} expresses the register-resident
analysis of Section~\ref{sec:method} directly: for each coefficient position the
kernel iterates the $k\times k$ tap grid with \verb|tl.static_range|, addresses
the corresponding $2\times2$ input block, forms the four subband values, and
accumulates into four \texttt{fp32} vectors. Because the kernel loads a compact
four-element weight vector per tap rather than an effective $(C,4,k,k)$ kernel
per output, its weight traffic is a quarter of the naive fused form. Block size,
warp count and pipeline depth are left to autotuning over eleven configurations
keyed on problem size.

The backward pass mirrors the forward: it iterates the same grid, applies the
transposed butterfly to the incoming subband gradients, and writes to a
\emph{disjoint} $2\times2$ input block per thread, so no atomics are required for
$\nabla_{\!X}$. Weight and scale gradients use Eq.~\eqref{eq:foldgrad} over
recomputed coefficients.

Performance matches the CUDA realisation closely --- slightly ahead in
\texttt{fp32} ($2.71$--$3.09\times$ against $2.24$--$2.89\times$) and slightly
behind in \texttt{fp16} --- which is the expected outcome if both are limited by
the same traffic. Two practical differences are worth recording. Triton's JIT
removes the dependency on a local \texttt{nvcc} installation, and its autotuner
adapts block sizes to the device without hand-tuning. Against that, the kernel
specialises on tensor shape, so a workload with many distinct shapes pays
repeated compilation, and the layer opts out of Dynamo tracing, so it cannot
currently be captured by \texttt{torch.compile}.

\subsection{Apple Metal}

The Metal realisation targets Apple Silicon through PyTorch's MPS backend. The
forward cascade uses a $16\times16$ threadgroup per $32\times32$ input tile,
staging the low-pass band in a single threadgroup buffer that is aliased and
overwritten in place at each level ($256\to64\to16\to4$ floats) with threads
retiring geometrically between barriers. The inverse cascade follows
Proposition~\ref{prop:bitindexed} per thread, exactly as in CUDA. As on NVIDIA,
\texttt{fp16} kernels convert to \texttt{fp32} on load, keep all butterflies and
all cross-level accumulation in \texttt{fp32}, and round once at the store; the
inter-level threadgroup buffer is \texttt{fp32} regardless of input precision.

Measured speedups are $2.05$--$2.40\times$. Three caveats attach to these
numbers and we would not want them read as a vendor comparison. First, the
baseline is PyTorch's generic MPS backend, which is substantially less mature
than cuDNN; this inflates the Apple-side ratios, and the dense-convolution
column in particular ($5\times$--$13\times$ slower than the fused layer) reflects
the maturity of the baseline far more than it reflects our kernels. Second, the
current dispatch issues a full stream synchronisation before encoding each
operation, which serialises host and device and makes these figures pessimistic
for end-to-end throughput on both arms. Third, several inverse dispatches
currently launch single-thread threadgroups, so occupancy is far below what the
formulation permits --- the Metal numbers should be read as a lower bound on
what the reformulation can deliver on this hardware rather than as a tuned
result.

\subsection{Coverage}

The CUDA and Triton paths support \texttt{fp32}, \texttt{fp16} and
\texttt{bf16}; the Metal path supports \texttt{fp32} and \texttt{fp16} only. All
three implement Haar with $L\le5$, depthwise, $C_{\mathrm{in}}=C_{\mathrm{out}}$,
and odd $k$.
